\chapter{Conclusiones}
\label{ch:conclusiones}

Se cumplió el objetivo del proyecto,
que consistía en diseñar un prototipo de sistema automático para la detección y
reconocimiento de caracteres en placas vehiculares costarricenses 
mediante técnicas de visión por computadora y aprendizaje profundo.
Además, se logró establecer una arquitectura basada de dos etapas basadas en aprendizaje profundo
---YOLOv11 para la detección de placas vehiculares y Fast-Plate-OCR para el reconocimiento óptico de caracteres---
evitando así las limitaciones de los sistemas tradicionales basados en segmentación de caracteres. 
La integración en un único flujo de procesamiento permitió validar el correcto funcionamiento del sistema completo,
desde la detección de las placas hasta el reconocimiento de texto correspondiente.

En cuanto a la preparación del conjunto de datos,
se generaron más de 8000 imágenes sintéticas,
junto con un esquema de aumento de datos con los cuales entrenar el 
módulo de reconocimiento de caracteres. 
Esta estrategia permitió ampliar la variabilidad de entrenamiento sin depender de
exclusivamente de la recolección de imágenes reales.

Los mejores resultados experimentales en la detección de placas se consiguieron 
al enriquecer el conjunto de datos de entrenamiento con placas vehiculares de otro país,
esto le permite al modelo generalizar mejor el concepto de ``placa'' y evita el 
sobreajuste a las imágenes de vehículos costarricenses.
De manera análoga, el mejor resultado experimental en el reconocimiento de placas
vehiculares se obtuvo al generar aumentos de datos con distintos tipos de aumento,
al eliminar tipos de aumento individualmente, la capacidad de generalización
del sistema decae.

Respecto a la evaluación del sistema,
se establecieron métricas de desempeño y se comparó el sistema propuesto con modelos base.
Los resultados experimentales demuestran un desempeño sobresaliente en ambos módulos principales. 
El detector YOLO alcanzó una precisión promedio del 93.5\%, y un mAP@0.5 de 0.946,
superando holgadamente el criterio de mejora mínima del 3\%, respecto al modelo base. 
Por su parte, el reconocedor de placas logró una precisión del 98.4\%,
superando en más de 90 puntos porcentuales al modelo base (3.7\,\%). 
Estos resultados evidencian la efectividad del entrenamiento con datos aumentados y el diseño modular del sistema.

En cuanto al desempeño del sistema integrado,
las pruebas sobre secuencias de video mostraron una latencia promedio de 300~ms por cuadro y
una velocidad de 3.84~FPS en una máquina sin aceleración por GPU. 
Si bien estas métricas no corresponden a operación en tiempo real de alta velocidad,
resultan adecuadas para el contexto del estacionamiento del Campus,
donde los vehículos se desplazan a baja velocidad y 
permanecen dentro del campo visual el tiempo suficiente para ser procesados.
Se espera que el uso de modelos optimizados (\texttt{ONNX}) y
hardware con aceleración CUDA incremente significativamente la velocidad del sistema, manteniendo su precisión.

En términos generales, el desarrollo permitió cumplir los objetivos planteados:
diseñar, entrenar e integrar un sistema robusto de detección y
reconocimiento de placas vehiculares, superando ampliamente las métricas base en precisión y exactitud completa.
El acceso al código fuente de este proyecto se encuentra en el siguiente
\href{https://github.com/jcorea510/TFG-ANPR.git}{\textit{enlace}} a un repositorio de control de versiones.

\section{Recomendaciones}
Este sistema presenta limitaciones que,
si se desea utilizar en la implementación posterior del sistema gestión vehicular, es recomendable abordar.
Con la escasa cantidad de datos disponibles, la validación del modelo 
de detección de placas puede estar sesgada, haciendo necesario incrementar
la cantidad de datos disponibles; y en la misma línea, en 
la sección de evaluación de fallos se encontró que el desempeño
de los sistemas de detección y reconocimiento de placas
se encuentra afectado debido a errores en las etiquetas de los datos,
por lo que conviene revisar las etiquetas del conjunto de datos empleado.
Por otro lado, en lo que respecta al reconocimiento de placas, 
en el conjunto de datos no se encuentran todos los tipos de placas vehiculares de Costa Rica.
Esto impone limitantes a los escenarios donde el sistema puede ser utilizado.
Es necesario, por lo tanto, no solo enriquecer el conjunto de datos
con automóviles que tengan más tipos de placas diferentes, sino 
también hacer más robusto el programa de generación sintética de placas
de entrenamiento, para que sea capaz de generar esas placas.
En adición, conviene realizar una mayor exploración en el aumento de datos, 
de manera que se puedan tratar los casos de fallos encontrados en los experimentos.
Finalmente, en la integración del sistema completo, resulta conveniente
evaluar el rendimiento con el sistema optimizado y en el hardware final 
donde se ejecutará la aplicación.

En síntesis, el sistema propuesto constituye una base sólida para aplicaciones 
prácticas de reconocimiento automático de placas vehiculares,
demostrando mejoras sustanciales respecto a modelos previos y
abriendo la posibilidad de futuras extensiones hacia escenarios en tiempo real o
sistemas inteligentes de gestión vehicular.

